% Paper 1, §8 — Normalization axis
% Thesis: LayerNorm → RMSNorm; Post-LN → Pre-LN (with 1/√N residual init)
% → DeepNorm. Why frontier-2026 is RMSNorm + Pre-LN with 1/√N init.

\section{Normalization}
\label{sec:normalization-axis}
\label{sec:norm}

The \glsterm{transformer}{Transformer}'s normalization layer evolved along two orthogonal axes
between 2017 and 2026: the \emph{type} of normalization (\glsterm{layer-normalization}{LayerNorm} vs
\glsterm{rmsnorm}{RMSNorm}) and the \emph{placement} of the norm relative to the residual
addition (\glsterm{post-ln}{Post-LN} vs \glsterm{pre-ln}{Pre-LN}, plus stabilizer variants such as DeepNorm).
The frontier-2026 consensus across LLaMA, \glsterm{mistral}{Mistral}, DeepSeek, and Qwen is
\textbf{\glsterm{rmsnorm}{RMSNorm} + \glsterm{pre-ln}{Pre-LN} with $1/\sqrt{N}$ residual scaling at
initialization}; the \glsterm{parameters}{parameters} of this consensus are the smallest
non-trivial slice of the architecture, and each piece has a directly
attributable primary source.

\subsection{LayerNorm: the canonical 2016 definition}
\label{sec:norm-layernorm}

\glsterm{layer-normalization}{LayerNorm} normalizes across the \glsterm{feature}{feature} axis of a single example, so
that the same statistics are computed at training and inference and
batch size has no effect. For an input vector
$\mathbf{a}^l \in \mathbb{R}^H$ (the summed inputs to the $H$ hidden
units of layer $l$),~\citet{ba2016-layernorm} define%
\footnote{KB excerpt: \texttt{kb/excerpts/ba2016-layernorm.md\#sec-3}.}
\begin{equation}
  \mu^l = \frac{1}{H} \sum_{i=1}^{H} a_i^l,
  \qquad
  \sigma^l = \sqrt{\frac{1}{H} \sum_{i=1}^{H} (a_i^l - \mu^l)^2},
  \tag{Ba--Eq.~3}
  \label{eq:ba-eq3}
\end{equation}
and apply a per-\glsterm{feature}{feature} gain $\mathbf{g}$ and bias $\mathbf{b}$ after
the normalization but before the non-linearity:
\begin{equation}
  \bar{\mathbf{a}}^l \;=\; \frac{\mathbf{g}}{\sigma^l} \odot
  (\mathbf{a}^l - \mu^l) + \mathbf{b}.
  \label{eq:ba-eq3-affine}
\end{equation}
Two structural properties follow directly from
Eq.~\eqref{eq:ba-eq3}--\eqref{eq:ba-eq3-affine}: (i) all hidden units in
a layer share the same $\mu^l, \sigma^l$ but \emph{different training
cases} get different normalization terms~\citep[\S 3]{ba2016-layernorm};
(ii) \glsterm{layer-normalization}{LayerNorm} is invariant under re-scaling of the entire weight
matrix and under re-centering of the incoming
weights~\citep[\S 4, Table~1]{ba2016-layernorm}. The second invariance
is what \citet{zhang2019} will later argue is dispensable.

In modern \glsterm{decoder-only}{decoder-only} notation, $H$ is the model dimension $\mathsf{D}$
and $\mathbf{a}^l$ is a single token's \glsterm{residual-stream}{residual stream} slice at depth
$l$; the normalization is applied independently per token, per layer.
The shape passing through is $\BSH$.

\subsection{RMSNorm: dropping the mean}
\label{sec:norm-rmsnorm}

\citet{zhang2019} hypothesize that \glsterm{layer-normalization}{LayerNorm}'s success rests on its
\emph{re-scaling} invariance, not its \emph{re-centering} invariance,
and propose Root Mean Square \glsterm{layer-normalization}{Layer Normalization}%
\footnote{KB excerpt: \texttt{kb/excerpts/zhang2019.md\#sec-3-1}.}
defined verbatim as
\begin{equation}
  \bar{a}_i \;=\; \frac{a_i}{\mathrm{RMS}(\mathbf{a})}\, g_i,
  \qquad
  \mathrm{RMS}(\mathbf{a}) \;=\;
  \sqrt{\frac{1}{n} \sum_{i=1}^{n} a_i^2}.
  \tag{Zhang--Eq.~4}
  \label{eq:rmsnorm}
\end{equation}
Comparing Eq.~\eqref{eq:rmsnorm} to Eq.~\eqref{eq:ba-eq3}: \glsterm{rmsnorm}{RMSNorm} drops
the mean $\mu^l$ entirely (and with it the bias $\mathbf{b}$), keeping
only the scale parameter $\mathbf{g}$. Two reductions over the \glsterm{feature}{feature}
axis become one; one element-wise add disappears.

Empirically, dropping the mean is ``free.'' \citet{zhang2019} report
matched quality across machine translation, reading comprehension,
summarization, and image classification, with runtime savings of
\textbf{7\%--64\%} depending on architecture and norm
placement~\citep[\S 4]{zhang2019}; on \glsterm{transformer}{Transformer}-base WMT'14 EN--DE
the swap costs $0.05$ BLEU and saves 9.4\% of training
time~\citep[\S 4]{zhang2019}. The shift-invariance that \glsterm{layer-normalization}{LayerNorm}
afforded turns out not to be load-bearing in
practice~\citep[\S 3.2]{zhang2019}, and every frontier-2026 \glsterm{decoder-only}{decoder-only}
\glsterm{llm}{LLM} (LLaMA-1/2/3, \glsterm{mistral}{Mistral}, Qwen, DeepSeek-V2/V3, Gemma, OLMo, Phi-4)
ships \glsterm{rmsnorm}{RMSNorm} rather than
\glsterm{layer-normalization}{LayerNorm}~\citep{touvron2023,jiang2023,meta-llama3,deepseek-v3,qwen3,olmo2,gemma3}.

\subsection{Post-LN versus Pre-LN}
\label{sec:norm-placement}

Independent of \emph{which} norm you compute, the question of
\emph{where} you compute it is a separate axis. Writing $F$ for a
sublayer (attention or \glsterm{ffn}{FFN}) and $\mathbf{x}$ for its input, the
canonical placements are
\begin{align}
  \text{Post-LN} \quad & \mathbf{y} \;=\; \mathrm{Norm}(\mathbf{x} + F(\mathbf{x}))
  \label{eq:post-ln} \\
  \text{Pre-LN} \quad & \mathbf{y} \;=\; \mathbf{x} + F(\mathrm{Norm}(\mathbf{x})).
  \label{eq:pre-ln}
\end{align}
\citet{vaswani2017} used \glsterm{post-ln}{Post-LN} \eqref{eq:post-ln} in the original
\glsterm{transformer}{Transformer}; the norm sits \emph{on} the residual path, so the residual
stream itself is renormalized after every block.~\citet{radford2019-gpt2}
moved the norm in GPT-2, in a passage worth quoting
verbatim%
\footnote{KB excerpt: \texttt{kb/excerpts/radford2019-gpt2.md\#sec-2-3}.}:
\begin{quote}
``Layer normalization (Ba et al., 2016) was moved to the input of each
sub-block, similar to a pre-activation residual network (He et al.,
2016) and an additional layer normalization was added after the final
\glsterm{self-attention}{self-attention} block. A modified initialization which accounts for the
accumulation on the residual path with model depth is used. We scale
the weights of residual layers at initialization by a factor of
$1/\sqrt{N}$ where $N$ is the number of residual layers.''
\citep[\S 2.3]{radford2019-gpt2}
\end{quote}
Two architectural decisions are bundled in that paragraph: the \glsterm{pre-ln}{Pre-LN}
move \eqref{eq:pre-ln}, and the $1/\sqrt{N}$ residual-path scaling at
initialization. The first changes the placement; the second compensates
for the fact that an unnormalized \glsterm{residual-stream}{residual stream}'s variance grows with
depth. Both are still in effect in frontier-2026 models, with $N$
typically the count of attention~+~\glsterm{ffn}{FFN} sublayers (i.e., $2L$ for an
$L$-block \glsterm{decoder-only}{decoder-only} stack).

\subsection{The Pre-LN stability proof}
\label{sec:norm-stability}

\citet{xiong2020} furnish the theoretical justification for the GPT-2
move via a mean-field analysis of gradient norms at
initialization%
\footnote{KB excerpt: \texttt{kb/excerpts/xiong2020.md\#sec-4}.}.
Stating their two main results:
\begin{theorem}[\glsterm{post-ln}{Post-LN} gradient scaling at init,
\citealp{xiong2020}, Thm.~1]
\label{thm:postln}
For the \glsterm{post-ln}{Post-LN} \glsterm{transformer}{Transformer} with standard Xavier initialization, the
expected Frobenius gradient norm at the final \glsterm{ffn}{FFN} sublayer is
\[
  \|\partial \mathcal{L} / \partial W^{2,L}\|_F
  \;=\; O\!\left( d \sqrt{\ln d} \right),
\]
independent of depth $L$. The gradient propagated through the residual
stream to earlier layers accumulates an additional $O(\sqrt{L})$ near
the output: parameter updates near the output layer become
disproportionately large as $L$ grows.
\end{theorem}
\begin{theorem}[\glsterm{pre-ln}{Pre-LN} gradient scaling at init,
\citealp{xiong2020}, Thm.~2]
\label{thm:preln}
For the \glsterm{pre-ln}{Pre-LN} \glsterm{transformer}{Transformer} with the same initialization, the expected
gradient norm at every layer is
\[
  \|\partial \mathcal{L} / \partial W^{l}\|_F
  \;=\; O\!\left( d \sqrt{\ln d / L} \right),
\]
balanced across depth and bounded as $L \to \infty$.
\end{theorem}
The practical consequence~\citep[\S 4.1, \S 5]{xiong2020}: \glsterm{post-ln}{Post-LN}
training requires the linear-\glsterm{warmup}{warmup} schedule
$\mathrm{lr}(t) = (t/T_{\mathrm{warmup}})\,\mathrm{lr}_{\max}$ to
compensate for the output-layer gradient blow-up, while \glsterm{pre-ln}{Pre-LN} trains
stably \emph{without \glsterm{warmup}{warmup}} at the same maximum learning rate. Their
WMT'14 and BERT-pretraining experiments confirm the prediction:
\glsterm{post-ln}{Post-LN} with $\mathrm{lr}_{\max} = 10^{-3}$ and no \glsterm{warmup}{warmup} diverges, but
the same setting with $T_{\mathrm{warmup}} = 4000$ trains stably; the
\glsterm{pre-ln}{Pre-LN} counterpart trains stably at $\mathrm{lr}_{\max} = 10^{-3}$ with
no \glsterm{warmup}{warmup} at all~\citep[\S 5]{xiong2020}. Theorem~\ref{thm:preln} is the
load-bearing reason every 2020+ \glsterm{decoder-only}{decoder-only} \glsterm{llm}{LLM} is \glsterm{pre-ln}{Pre-LN}.

\subsection{DeepNorm and beyond}
\label{sec:norm-deepnorm}

A natural follow-up question is whether one could keep the \glsterm{post-ln}{Post-LN}
placement (which has its own expressivity arguments at very deep stacks)
but rescale the residual contribution to defang
Theorem~\ref{thm:postln}. \citet{wang2022-deepnorm}'s DeepNet answers
yes: introduce a per-block constant $\alpha > 1$ on the residual
addition,
\begin{equation}
  \mathbf{y} \;=\; \mathrm{Norm}(\alpha \mathbf{x} + F(\mathbf{x})),
  \label{eq:deepnorm}
\end{equation}
with $\alpha$ and the sublayer initialization scale chosen as functions
of depth $L$ such that the \glsterm{residual-stream}{residual stream}'s variance and the
gradient's expected norm both stay $O(1)$ at initialization for $L$ as
large as $1000$.~\deepencite{wang2022-deepnorm \S 3 (DeepNorm scaling
constants $\alpha, \beta$ as functions of $L$ for encoder-only,
decoder-only, encoder-decoder)} DeepNorm thus generalizes the \glsterm{pre-ln}{Pre-LN} /
$1/\sqrt{N}$ trick to a regime where \glsterm{pre-ln}{Pre-LN}'s residual-stream-norm
drift becomes a problem on its own, and demonstrates that the
1000-layer regime is reachable. It is used in some Microsoft tech-report
models but has not displaced \glsterm{pre-ln}{Pre-LN} at frontier scale; modern frontier
LLMs run at $L = 32\text{--}126$ layers~\citep{deepseek-v3,meta-llama3}
where \glsterm{pre-ln}{Pre-LN} is comfortable and the DeepNorm machinery is not yet
necessary.

\begin{intuition}
\glsterm{layer-normalization}{LayerNorm} is ``make every token's \glsterm{residual-stream}{residual stream} have the same energy
along the \glsterm{feature}{feature} axis, then re-scale per-channel.'' Energy of
$\mathbf{a}$ along the \glsterm{feature}{feature} axis is $\sum_i a_i^2$, so \glsterm{rmsnorm}{RMSNorm}
(Eq.~\eqref{eq:rmsnorm}) is the same statement with the offset removed:
divide each component by $\sqrt{(1/n) \sum_i a_i^2}$ so the resulting
vector has unit RMS, then apply a learned per-channel gain $g_i$. The
mean-subtraction in \glsterm{layer-normalization}{LayerNorm} pre-projects onto the
$\mathbf{1}^\perp$-hyperplane; if the model doesn't use that
projection's invariance (Zhang \& Sennrich's \S 3.2 finding), dropping
it costs nothing. Returning to math: \glsterm{rmsnorm}{RMSNorm} preserves
$\mathrm{RMS}(\bar{\mathbf{a}}) = \|\mathbf{g}\|_2 / \sqrt{n}$; the
energy of the normalized stream is fixed by $\mathbf{g}$ alone.
\end{intuition}

\subsection{Why frontier-2026 settled here}
\label{sec:norm-frontier}

Composing the four moves: take \glsterm{layer-normalization}{LayerNorm}
(Eq.~\eqref{eq:ba-eq3}--\eqref{eq:ba-eq3-affine}), drop the mean and
the bias to get \glsterm{rmsnorm}{RMSNorm} (Eq.~\eqref{eq:rmsnorm}); place the norm
\emph{before} each sublayer per Eq.~\eqref{eq:pre-ln}; scale the
residual-path weights at init by $1/\sqrt{N}$ per
\citet[\S 2.3]{radford2019-gpt2}; rely on Theorem~\ref{thm:preln} for
\glsterm{warmup}{warmup}-free trainability. The total cost on top of the canonical 2017
\glsterm{transformer}{Transformer} is one bias parameter removed, one mean-subtraction
removed, one norm relocated, and one initialization scale changed --
each backed by a primary source. \citet{touvron2023}'s LLaMA-1
specification adopts the full bundle, and every subsequent
LLaMA-/\glsterm{mistral}{Mistral}-/DeepSeek-/Qwen-class model has shipped some variant
of it~\citep{jiang2023,meta-llama3,deepseek-v3,qwen3}, with two open
splits: 2025+ models add QK-Norm to the attention path
(\citealp{olmo2}; covered in \S\ref{sec:attention-mechanism-axis} of
this paper), and a separate stream switches to Peri-LN for very-deep
stacks~\citep{olmo2,gemma3,peri-ln2025}. Both are extensions of, not
replacements for, the \glsterm{rmsnorm}{RMSNorm} + \glsterm{pre-ln}{Pre-LN} + $1/\sqrt{N}$ baseline
established between 2016 and 2020.

The next section (\S\ref{sec:state-space-and-hybrids}) leaves the
attention-and-\glsterm{ffn}{FFN}-with-residuals architecture entirely and asks what
happens when the \glsterm{residual-stream}{residual stream} is replaced by a state-space
recurrence -- a setting in which the normalization choices of this
section largely transfer (state-space-model-based LLMs use \glsterm{rmsnorm}{RMSNorm} and
analogues of \glsterm{pre-ln}{Pre-LN}), but the gradient-scaling argument of
Theorem~\ref{thm:preln} no longer applies because the residual
structure is gone.
